\documentclass[11pt,a4paper]{article}
\usepackage{operationsform}
\begin{document}
\OPStitle{取材質問設計表}{TPL-DOC-046}{0.2.0}
\OPSsection{1}{基本情報}
\OPSblanklong{案件・文書の識別子}
\OPSblankshort{作成者と連絡先}
\OPSblankshort{作成日と対象版}
\OPSsection{2}{質問の設計}
\OPSblanklong{取材目的と中心となる問い}
\OPSblanklong{対象者の専門的背景}
\OPSblanklong{既知の事実と確認すべき不確実性}
\OPSblanklong{質問設計上の倫理}
\OPSsection{3}{質問と根拠}
\par\Needspace{4\baselineskip}\noindent\textbf{質問・意図・根拠資料・追質問}\par\vspace{1mm}
{\fontsize{8.5}{13}\selectfont
\begin{tabularx}{\linewidth}{@{}XXXX@{}}
\textbf{質問ID・本文} & \textbf{意図と根拠ID} & \textbf{想定追質問} & \textbf{時間・優先順位}\\[2mm]
\strut & \strut & \strut & \strut\\[5mm]
\strut & \strut & \strut & \strut\\[5mm]
\strut & \strut & \strut & \strut\\[5mm]
\end{tabularx}
}\vspace{2mm}\par
\OPSblanklong{優先順位と所要時間}
\OPSblanklong{予想される難しい論点}
\OPSblanklong{発言確認の方法}
\OPSsection{4}{確認と履歴}
\par\Needspace{4\baselineskip}\noindent\textbf{参照資料と根拠ID}\par\vspace{1mm}
{\fontsize{8.5}{13}\selectfont
\begin{tabularx}{\linewidth}{@{}XXXX@{}}
\textbf{資料ID・箇所} & \textbf{支持する事実} & \textbf{未確認の論点} & \textbf{確認担当}\\[2mm]
\strut & \strut & \strut & \strut\\[5mm]
\strut & \strut & \strut & \strut\\[5mm]
\strut & \strut & \strut & \strut\\[5mm]
\end{tabularx}
}\vspace{2mm}\par
\OPSblanklong{未確認事項・例外}
\par\Needspace{4\baselineskip}\noindent\textbf{点検・改訂の記録}\par\vspace{1mm}
{\fontsize{8.5}{13}\selectfont
\begin{tabularx}{\linewidth}{@{}XXXX@{}}
\textbf{対象・記録ID} & \textbf{内容・根拠} & \textbf{確認者・日付} & \textbf{状態・備考}\\[2mm]
\strut & \strut & \strut & \strut\\[5mm]
\strut & \strut & \strut & \strut\\[5mm]
\strut & \strut & \strut & \strut\\[5mm]
\end{tabularx}
}\vspace{2mm}\par
\end{document}
